\section{Native support-frame coefficients and actual local section values}

\begin{lemma}\label{lem:fc-support-factor-scalar-residue}
\lean{EnvelopeProperty.Berkovich.SchemeSupportFrameEvaluation.factorFrameCoefficient,EnvelopeProperty.Berkovich.SchemeSupportFrameEvaluation.supportFactor_appTop}
\leanok
\uses{def:scheme-support-line-frame,lem:fc-native-unit-section-morphism,lem:fc-original-line-section-frame-naturality,scheme:actual-residue-section-pullback}
Let $y$ be any actual native valued point of a scheme $Y$, and
let $V$ be an original open containing its support. For the actual
residue factor $q:\operatorname{Spec}\kappa_Y(y)\to V$, applying
$q$ to a whole scalar section of $V$ and then using the genuine
affine section isomorphism gives exactly the original residue
evaluation of the corresponding section on $V\subseteq Y$.
Moreover, for any actual factor $q b=h$ and actual original frame
on $b^*N$, its transported frame applied to the native pullback
of every original whole section is exactly the actual scheme
pullback of that original frame's whole scalar coefficient.
\end{lemma}
\begin{proof}
For the first identity, the actual whole triangle $q\iota=h_Y$
gives the composition identity for local section maps. The open
inclusion map identifies sections on $V$ with whole sections of
its open subscheme. Following the composition identity by the
genuine affine section isomorphism is the actual original residue
evaluation formula. For the second identity, express each native
pullback-unit section as the value at one of the actual pulled
section morphism. Applying the given frame produces the actual
scalar endomorphism of its value on one. The native factor-frame
naturality theorem pulls this entire endomorphism through $q$.
Genuine scalar pullback sends it to multiplication by the actual
scheme pullback of its original scalar. Evaluate at one to obtain
the asserted whole section coefficient identity.
\end{proof}

\begin{theorem}\label{thm:fc-native-support-frame-section-evaluation}
\lean{EnvelopeProperty.Berkovich.SchemeSupportFrameEvaluation.localCoefficient,EnvelopeProperty.Berkovich.SchemeSupportFrameEvaluation.supportFrame_coefficient,EnvelopeProperty.Berkovich.SchemeSupportFrameEvaluation.supportFrame_value}
\leanok
\uses{lem:fc-support-factor-scalar-residue,def:scheme-support-line-frame}
Let $N$ be any original module on $Y$ and let
$e_V:N|_V\simeq\mathcal O_V$ be an actual entire original frame.
For every original whole section $\sigma$ of $N$, restrict it by
the actual native pullback unit to the open subscheme and apply
$e_V$. Regard the resulting scalar as the actual original section
$a_\sigma$ of $\mathcal O_Y$ on $V$. Its native residue evaluation
at $y$ is exactly the coefficient of $\sigma$ in the genuine
support-induced residue frame. Consequently the actual native
valuation of that coefficient equals the actual native section
value $|a_\sigma(y)|$.
\end{theorem}
\begin{proof}
The support-induced residue frame is the actual factor frame
through $q$. The preceding factor-frame formula identifies its
coefficient before the affine section isomorphism with the true
pullback through $q$ of the actual open coefficient. The preceding
scalar residue formula identifies the resulting scalar with the
actual residue evaluation of $a_\sigma$. Apply the actual native
residue valuation at $y$. By definition this is native section
evaluation on $V$. Every original section is included, also zero;
no properness, bounded base, frame matching or norm input occurs.
\end{proof}

\begin{lemma}\label{lem:fc-native-open-frame-lift-evaluation}
\lean{EnvelopeProperty.Berkovich.SchemeSupportFrameEvaluation.eval_openTopIso,EnvelopeProperty.Berkovich.SchemeSupportFrameEvaluation.supportFrame_value_of_openLift}
\leanok
\uses{thm:fc-native-support-frame-section-evaluation,scheme:valued-open-immersion-image}
Let $z$ be any actual native point of the open subscheme $V$.
The native value at its image in $Y$ of an original scalar section
on $V$ equals its whole-section value at $z$. In particular, the
support-frame coefficient value of an original whole section
$\sigma$ at this image equals the value at $z$ of its actual
whole scalar coefficient in $e_V$.
\end{lemma}
\begin{proof}
The actual inverse section isomorphism for the open immersion
preserves native evaluation. On the whole open subscheme its
image open is $V$ and that genuine inverse section map is exactly
the original open whole-section identification. This proves the
first identity. Apply the preceding support-frame theorem to the
image point and the actual coefficient section. The open
whole-section identification and its inverse cancel, leaving
exactly the true whole scalar coefficient at $z$.
\end{proof}

\begin{corollary}\label{cor:fc-native-laurent-support-frame-values}
\lean{EnvelopeProperty.Berkovich.SchemeSupportFrameEvaluation.laurentSupportFrame,EnvelopeProperty.Berkovich.SchemeSupportFrameEvaluation.laurentSupportFrame_value,EnvelopeProperty.Berkovich.SchemeSupportFrameEvaluation.exists_laurentSupportFrame_values,EnvelopeProperty.Berkovich.SchemeSupportFrameEvaluation.exists_laurentSupportFrame_values_of_isInvertible,EnvelopeProperty.Berkovich.SchemeSupportFrameEvaluation.laurentSupportFrame_finite_family}
\leanok
\uses{lem:fc-native-open-frame-lift-evaluation,thm:fc-point-laurent-local-coefficient-value}
For an original invertible $M$ on a proper compact separated
$X\to\operatorname{Spec}R$, an actual bounded-base relative point
$x$, any actual frame on a neighborhood $U$ of its true proper
center, and every positive radius $r$, the actual product Gauss
point lies in the original product frame open. The actual product
frame there induces an entire original residue frame $e_0$ such
that for every original whole Laurent-product section $\sigma$,
\[
 |c_{e_0}(\sigma)|=\max_i\|m_i\|_{M,x}r^i.
\]
Here the $m_i$ are its actual original Laurent coefficients and
the norms are the true original proper-model line norms of their
native residue-unit sections. One such entire residue frame works
for every original section. Original invertibility itself supplies
such a frame, without an open or local trivialization as an input. For every finite original family the
same original finite coefficient and integer-weight family,
chosen before all points and radii, computes its maximum.
\end{corollary}
\begin{proof}
The actual product-open lift of the source Gauss point maps as a
whole native valued point to the original product Gauss point.
In particular its support is in the product frame open. Apply
the preceding lemma to this actual lift and the genuine product
local frame. Its true whole scalar coefficient is exactly the
original product local coefficient. The native Laurent local
coefficient theorem computes its value by the stated original
weighted maximum, with every negative exponent and zero included.
The entire support-induced frame has no section input, so it
works for all original sections. For the version with no chosen local
frame, the genuine invertible open cover contains the true proper
center in one member; its actual frame supplies the preceding
construction. Taking finite maxima and using
the original finite Laurent coefficient-family theorem gives the
same coefficient and integer-weight family before all points and
radii. No frame, evaluation, point or norm comparison is assumed.
\end{proof}
